% A RESTORE NAMES THINGS BEHIND THE ESCAPE CHARACTER, whatever
% \escapechar says that is -- and nothing at all where it is out of
% range:
%
%   \escapechar=`|     {restoring |count5=7}   {restoring |m=macro:->a}
%   \escapechar=-1     {restoring count5=7}    {restoring m=macro:->a}
%   \escapechar=`\\    {restoring \count5=7}
%
% Every name in the line takes it: the register or parameter, a macro's
% own name, a family's font (`{restoring |textfont2=|ff}' has two), the
% font in use, and the `\ETC.' that ends a value shown to its limit.
% hstex wrote a backslash into all of them.
%
% What is NOT behind it: `current font=' is words rather than a name, and
% so are `macro:', `undefined' and `void'. NOR IS AN ACTIVE CHARACTER,
% which is written as itself whatever \escapechar says -- trip's
% `{\catcode`?=13 ...}' gives back `{restoring ?=undefined}'.
%
% RUN THIS UNDER `pdftex -ini'; it sets its own catcodes.
\catcode`\{=1 \catcode`\}=2 \catcode`\#=6
\tracingonline=1 \tracingrestores=1
\count5=7 \dimen5=1pt \skip5=1pt \toks5={a}\def\m{a}\catcode`\q=12
\fam=3 \everypar={p}\parshape=1 0pt 100pt \font\ff=cmr10 \textfont2=\ff
\message{[A escapechar is a bar]}
\escapechar=`|
{\count5=8 \dimen5=2pt \skip5=2pt \toks5={b}\def\m{b}\catcode`\q=13
 \fam=4 \everypar={q}\parshape=1 0pt 50pt \textfont2=\nullfont \undefinedxx}
\message{[B escapechar out of range]}
\escapechar=-1
{\count5=9 \def\m{c}\fam=5 \let\zz=\relax}
\message{[C escapechar back]}
\escapechar=`\\
{\count5=10 \def\m{d}}
\message{[D the font in use]}
\escapechar=`|
\ff
{\nullfont}
\escapechar=`\\
\message{[E active characters]}
{\catcode`?=13 \def?{x}}
{\catcode`!=13 \let!=\relax}
\escapechar=-1
{\catcode`;=13 \def;{y}}
\escapechar=`\\
\message{[done]}
\end
